\fontsize{9pt}{11.8pt}\selectfont

\begin{multicols}{2}

\marginicon\textbf{53.}\ Ước lượng tiệm cận ngang của hàm số
\[
f(x) = \frac{3x^3 + 500x^2}{x^3 + 500x^2 + 100x + 2000}
\]
bằng cách vẽ đồ thị $f$ với $-10 \le x \le 10$. Sau đó tìm phương trình của tiệm cận bằng cách tính giới hạn. Em giải thích sự sai khác này như thế nào?

\vspace{6pt}
\marginicon\textbf{54.}\ (a) Vẽ đồ thị hàm số
\[
f(x) = \frac{\sqrt{2x^2 + 1}}{3x - 5}
\]
Em quan sát thấy bao nhiêu tiệm cận ngang và tiệm cận đứng? Sử dụng đồ thị để ước lượng giá trị của các giới hạn
\[
\lim_{x \to \infty} \frac{\sqrt{2x^2 + 1}}{3x - 5} \qquad\text{và}\qquad \lim_{x \to -\infty} \frac{\sqrt{2x^2 + 1}}{3x - 5}
\]
(b) Bằng cách tính các giá trị của $f(x)$, hãy đưa ra các ước lượng bằng số cho các giới hạn trong phần (a).\\
(c) Tính giá trị chính xác của các giới hạn trong phần (a). Em nhận được cùng một giá trị hay các giá trị khác nhau cho hai giới hạn này? [Từ câu trả lời cho phần (a), em có thể cần kiểm tra lại phép tính của mình đối với giới hạn thứ hai.]

\vspace{6pt}
\textbf{55.}\ Cho $P$ và $Q$ là các đa thức. Tìm
\[
\lim_{x \to \infty} \frac{P(x)}{Q(x)}
\]
nếu bậc của $P$ là (a) nhỏ hơn bậc của $Q$ và (b) lớn hơn bậc của $Q$.

\vspace{6pt}
\textbf{56.}\ Hãy vẽ phác thảo đường cong $y = x^n$ ($n$ là một số nguyên) cho năm trường hợp sau:
\begin{center}
\begin{tabular}{@{}p{0.48\linewidth}p{0.48\linewidth}@{}}
(i) $n = 0$ & (ii) $n > 0$, $n$ lẻ \\
(iii) $n > 0$, $n$ chẵn & (iv) $n < 0$, $n$ lẻ \\
(v) $n < 0$, $n$ chẵn &
\end{tabular}
\end{center}
Sau đó sử dụng các hình phác thảo này để tìm các giới hạn sau:
\begin{center}
\begin{tabular}{@{}p{0.48\linewidth}p{0.48\linewidth}@{}}
(a) $\lim\limits_{x \to 0^+} x^n$ & (b) $\lim\limits_{x \to 0^-} x^n$ \\[6pt]
(c) $\lim\limits_{x \to \infty} x^n$ & (d) $\lim\limits_{x \to -\infty} x^n$
\end{tabular}
\end{center}

\vspace{4pt}
\textbf{57.}\ Tìm công thức cho một hàm số $f$ thỏa mãn các điều kiện sau:
\[
\lim_{x \to \pm\infty} f(x) = 0, \qquad \lim_{x \to 0} f(x) = -\infty, \qquad f(2) = 0,
\]
\[
\lim_{x \to 3^-} f(x) = \infty, \qquad \lim_{x \to 3^+} f(x) = -\infty
\]

\vspace{4pt}
\textbf{58.}\ Tìm công thức cho một hàm số có các tiệm cận đứng $x = 1$ và $x = 3$ cùng tiệm cận ngang $y = 1$.

\vspace{6pt}
\textbf{59.}\ Một hàm số $f$ là tỉ số của các hàm bậc hai, có một tiệm cận đứng $x = 4$ và chỉ có một giao điểm với trục hoành là $x = 1$. Biết rằng $f$ có một điểm gián đoạn khử được tại $x = -1$ và $\lim\limits_{x \to -1} f(x) = 2$. Hãy tính:
\begin{center}
\begin{tabular}{@{}p{0.48\linewidth}p{0.48\linewidth}@{}}
(a) $f(0)$ & (b) $\lim\limits_{x \to \infty} f(x)$
\end{tabular}
\end{center}

\vfill\null
\columnbreak

{\bfseries\color{stewartblue}60--64}\quad Tìm các giới hạn khi $x \to \infty$ và khi $x \to -\infty$. Sử dụng thông tin này, cùng với các giao điểm với các trục tọa độ, để vẽ phác thảo đồ thị như trong Ví dụ 12.

\vspace{4pt}
\textbf{60.}\ $y = 2x^3 - x^4$ \hfill \textbf{61.}\ $y = x^4 - x^6$

\vspace{3pt}
\textbf{62.}\ $y = x^3(x + 2)^2(x - 1)$

\vspace{3pt}
\textbf{63.}\ $y = (3 - x)(1 + x)^2(1 - x)^4$

\vspace{3pt}
\textbf{64.}\ $y = x^2(x^2 - 1)^2(x + 2)$

\vspace{6pt}
\textbf{65.}\ (a) Sử dụng Định lý kẹp để tính $\lim\limits_{x \to \infty} \dfrac{\sin x}{x}$.\\
\marginicon(b) Vẽ đồ thị hàm số $f(x) = (\sin x)/x$. Đồ thị cắt đường tiệm cận bao nhiêu lần?

\vspace{6pt}
\marginicon\textbf{66.}\ \textbf{Hành vi ở vô cực của một hàm số}\quad Khi nói đến \emph{hành vi ở vô cực} của một hàm số, ta muốn nói đến hành vi của các giá trị của nó khi $x \to \infty$ và khi $x \to -\infty$.

(a) Hãy mô tả và so sánh hành vi ở vô cực của các hàm số
\[
P(x) = 3x^5 - 5x^3 + 2x \qquad Q(x) = 3x^5
\]
bằng cách vẽ đồ thị cả hai hàm số trong các khung nhìn $[-2, 2] \times [-2, 2]$ và $[-10, 10] \times [-10\,000, 10\,000]$.\\
(b) Hai hàm số được gọi là có \emph{cùng hành vi ở vô cực} nếu tỉ số của chúng tiến tới 1 khi $x \to \infty$. Chứng minh rằng $P$ và $Q$ có cùng hành vi ở vô cực.

\vspace{6pt}
\textbf{67.}\ Tìm $\lim\limits_{x \to \infty} f(x)$ nếu với mọi $x > 1$,
\[
\frac{10e^x - 21}{2e^x} < f(x) < \frac{5\sqrt{x}}{\sqrt{x - 1}}
\]

\vspace{4pt}
\textbf{68.}\ (a) Một bể chứa $5000\text{ L}$ nước tinh khiết. Nước muối chứa $30\text{ g}$ muối trên mỗi lít nước được bơm vào bể với tốc độ $25\text{ L/phút}$. Hãy chứng minh rằng nồng độ muối sau $t$ phút (tính bằng gam trên lít) là
\[
C(t) = \frac{30t}{200 + t}
\]
(b) Điều gì xảy ra với nồng độ khi $t \to \infty$?

\vspace{6pt}
\textbf{69.}\ Trong Chương 9, ta sẽ có thể chứng minh được rằng, dưới các giả thiết nhất định, vận tốc $v(t)$ của một giọt mưa đang rơi tại thời điểm $t$ là
\[
v(t) = v^*(1 - e^{-gt/v^*})
\]
trong đó $g$ là gia tốc trọng trường và $v^*$ là \emph{vận tốc giới hạn} của giọt mưa.\\
(a) Tìm $\lim\limits_{t \to \infty} v(t)$.\\
\marginicon(b) Vẽ đồ thị $v(t)$ nếu $v^* = 1\text{ m/s}$ và $g = 9{,}8\text{ m/s}^2$. Cần bao lâu để vận tốc của giọt mưa đạt $99\%$ vận tốc giới hạn của nó?

\vspace{6pt}
\marginicon\textbf{70.}\ (a) Bằng cách vẽ đồ thị $y = e^{-x/10}$ và $y = 0{,}1$ trên cùng một màn hình, hãy tìm xem cần làm cho $x$ lớn đến mức nào để $e^{-x/10} < 0{,}1$.\\
(b) Em có thể giải phần (a) mà không dùng đồ thị không?

\end{multicols}

\vfill
\noindent{\fontsize{6.5pt}{8pt}\selectfont \color{black!70}
Copyright 2021 Cengage Learning. All Rights Reserved. May not be copied, scanned, or duplicated, in whole or in part. WCN 02-200-203\\
Due to electronic rights, some third party content may be suppressed from the eBook and/or eChapter(s). Editorial review has deemed that any suppressed content does not materially affect the overall learning experience. Cengage Learning reserves the right to remove additional content at any time if subsequent rights restrictions require it.\par}
